\section{Discussion}
\label{sec:discussion}

\paragraph{Defaults and state responses require different transfer claims.}
The evidence supports recurring educational behavior in these deployments.
Default action profiles learned on one source set improve predictions on
another; revelation and elicitation gains also survive the tested expression
changes without refitting. The stronger claim about student-state responses
has a different outcome. Acknowledgement interactions help under the original
wording, including its contemporary repetition, but lose their increment under
explicit synonyms and implicit cues. A successful new-problem test therefore
cannot by itself validate a transportable student-state profile.

This does not imply that models ignore student emotion: acknowledgement rates
still change between cue poles. It means that the particular conditional
probabilities learned under the old wording do not improve forecasts under
these new inputs. The explicit arm loses relative advantage despite slightly lower absolute
conditional loss; the implicit arm has higher conditional loss.
Expression-specific calibration, the additional information
in implicit cues, and changes in what the acknowledgement event captures are
plausible explanations. The study does not identify an internal mechanism or
establish failure of every possible conditional predictor. Educational
personality is consequently useful here as a testable action profile with
specified transfer boundaries, rather than a psychological type.

\paragraph{Prompt dependence has more than one form.}
ASK and EXPLAIN instructions make demanded actions nearly or fully shared,
while other permissible actions remain unevenly expressed. Neutral role
paraphrases cannot be declared population-equivalent with sixteen sources.
The peer control adds another distinction: social context can change whether
the tutor acknowledges a fixed current-student state. Selected examples show
that such an event can occur while correctly distinguishing the peer's emotion;
its presence alone does not diagnose emotional misattribution. Role wording,
student expression, and instructional constraint thus belong in the empirical
description of a deployment's behavior.

\paragraph{Increment over existing personality and tutoring research.}
Beyond prior situated profiles and tutor-action comparisons
\citep{lee2025trait,liu2026bma,kucheria2025cima,lee2026tutorpersonas}, the linked
design distinguishes default transfer, expression-limited conditional transfer,
and instructional convergence. Unchanged forecasts and the contemporary control
locate the predictive boundary. Call volume and automated coding are not the
novelty; education-specific traits, internal persona mechanisms, and teaching
benefits remain unestablished.

\paragraph{What the archive adds and leaves unresolved.}
Archived benchmark replies extend the evidence beyond constructed student
messages: model differences predict actions under general tutoring, while
explicit probing instructions sharply reduce that increment. Neither a single
deployment nor one coder accounts for the scaffolding result. Yet the archive
does not support an instruction-invariant absolute action profile; directly
transferred probabilities perform poorly. This can follow from instruction-wide
rate shifts without any reversal of model rankings. The short benchmark
contexts, rare acknowledgement events, and incomplete historical attestation
limit the result. Tutor, grading, and test-taking records still cannot be pooled
as a million independent observations of tutor personality.
